% !TeX program = pdflatex
% Versioni 3 - IEEEtran Conference.
% Për Overleaf: Menu -> Compiler -> pdfLaTeX.
% Burimet kryesore të kontekstit: Strategjia e Kosovës për Qeverisje Elektronike
% 2023-2027, Agjenda Digjitale e Kosovës 2030 dhe studimi KIPRED/DCAF (2022).
\documentclass[conference]{IEEEtran}
\usepackage[T1]{fontenc}
\usepackage[utf8]{inputenc}
\usepackage[albanian]{babel}
\usepackage{newtxtext}
\usepackage{cite}
\usepackage{graphicx}
\usepackage{booktabs}
\usepackage{array}
\usepackage{url}
\usepackage{hyperref}
\hypersetup{hidelinks}
% babel i rivendos emërtimet në \begin{document}, prandaj \addto
\addto\captionsalbanian{%
  \renewcommand{\abstractname}{Abstrakti}%
  \renewcommand{\refname}{Referencat}%
  \renewcommand{\tablename}{Tabela}%
}
\renewcommand{\IEEEkeywordsname}{Fjalët kyçe}
% Ndarja e autorëve në dy rreshta sipas renditjes në dokumentin origjinal.
\makeatletter
\newcommand{\linebreakand}{%
  \end{@IEEEauthorhalign}%
  \hfill\mbox{}\par
  \vspace{0.9em}%
  \mbox{}\hfill\begin{@IEEEauthorhalign}%
}
\makeatother
\begin{document}
% ---------------------------------------------------------------------
% 1. TITULLI I PUNIMIT
% ---------------------------------------------------------------------
\title{Modelet, korniza e mundësive dhe sfidat e e-Qeverisjes\\
në Kosovë: Vlerësim i integruar dhe propozim\\
për përmirësimin e koordinimit ndërinstitucional}
% Blloqet individuale të autorëve në formatin IEEE Conference: 3 + 2.
\author{
\IEEEauthorblockN{Shefket Bylygbashi}
\IEEEauthorblockA{\textit{Inxhinieri Kompjuterike}\\
\textit{dhe Softuerike}\\
\textit{Fakulteti i Inxhinierisë Elektrike}\\
\textit{dhe Kompjuterike}\\
Prishtinë, Kosovë\\
{\textit{shefket.bylygbashi@student.uni-pr.edu}}}\\
\and
\IEEEauthorblockN{Rinor Ukshini}
\IEEEauthorblockA{\textit{Inxhinieri Kompjuterike}\\
\textit{dhe Softuerike}\\
\textit{Fakulteti i Inxhinierisë Elektrike}\\
\textit{dhe Kompjuterike}\\
Prishtinë, Kosovë\\
{\textit{rinor.ukshini1@student.uni-pr.edu}}}\\
\and
\IEEEauthorblockN{Arton Gërguri}
\IEEEauthorblockA{\textit{Inxhinieri Kompjuterike}\\
\textit{dhe Softuerike}\\
\textit{Fakulteti i Inxhinierisë Elektrike}\\
\textit{dhe Kompjuterike}\\
Prishtinë, Kosovë\\
{\textit{arton.gerguri@student.uni-pr.edu}}}\\
\linebreakand
\IEEEauthorblockN{Rina Bunjaku}
\IEEEauthorblockA{\textit{Inxhinieri Kompjuterike}\\
\textit{dhe Softuerike}\\
\textit{Fakulteti i Inxhinierisë Elektrike}\\
\textit{dhe Kompjuterike}\\
Prishtinë, Kosovë\\
{\textit{rina.bunjaku@student.uni-pr.edu}}}\\
\and
\IEEEauthorblockN{Besarta Berisha}
\IEEEauthorblockA{\textit{Inxhinieri Kompjuterike}\\
\textit{dhe Softuerike}\\
\textit{Fakulteti i Inxhinierisë Elektrike}\\
\textit{dhe Kompjuterike}\\
Prishtinë, Kosovë\\
{\textit{besarta.berisha1@student.uni-pr.edu}}}\\
}
\maketitle
% Abstrakti do të plotësohet pasi të përfundojë punimi.
% \begin{abstract}
% \end{abstract}
% ---------------------------------------------------------------------
% 2. FJALËT KYÇE
% ---------------------------------------------------------------------
\begin{IEEEkeywords}
e-Qeverisja, ndërveprueshmëria, koordinimi ndërinstitucional
\end{IEEEkeywords}
% ---------------------------------------------------------------------
% 3. HYRJE
% ---------------------------------------------------------------------
\section{Hyrje}
\label{sec:hyrje}
Transformimi digjital i sektorit publik nuk është më një projekt
teknologjik, por një reformë institucionale që prek mënyrën se si shteti
organizon kompetencat, proceset dhe marrëdhënien me qytetarin. Paradigma e
\textit{qeverisjes së epokës digjitale} kërkon riintegrimin e funksioneve të
fragmentuara nga reformat e mëparshme menaxheriale, organizimin e
shërbimeve rreth nevojave të përdoruesit dhe digjitalizimin e proceseve
administrative nga fillimi deri në fund. Megjithatë,
teknologjia nuk prodhon vetvetiu transformim: sipas kornizës së
\textit{zbatimit të teknologjisë} (\textit{technology enactment}), efektet e
TIK-ut ndërmjetësohen nga strukturat burokratike, rregullat dhe rrjetet
ekzistuese ndërinstitucionale, të cilat shpesh e riprodhojnë fragmentimin në
vend që ta eliminojnë. Kjo shpjegon pse shumica e
projekteve të e-Qeverisjes në vendet në zhvillim dhe në tranzicion dështojnë
tërësisht ose pjesërisht, kryesisht për shkak të hendekut ndërmjet
projektimit formal dhe realitetit institucional.

Kosova përbën një rast studimi veçanërisht të pasur për këtë problematikë.
Vendi ka ndërtuar një trajektore të qëndrueshme strategjike: nga Strategjia
e parë e Qeverisjes Elektronike 2009--2015 dhe Agjenda Digjitale 2013--2020,
te anëtarësimi në Agjendën Digjitale të BE-së për Ballkanin Perëndimor
(2018) dhe Strategjia e Reformës së Administratës Publike (RAP) 2022--2027
\cite{strategjia2023}. Kjo trajektore kulmon me Agjendën Digjitale të
Kosovës 2030 (AD~2030), një agjendë horizontale ndërsektoriale me pesë
objektiva strategjikë dhe buxhet indikativ prej 152 milionë eurosh
\cite{agjenda2030}, si dhe me Strategjinë e Kosovës për Qeverisje
Elektronike 2023--2027, të hartuar me mbështetjen e Akademisë Estoneze të
e-Qeverisjes dhe të strukturuar në gjashtë objektiva: koordinimi e
menaxhimi, kompetencat digjitale, ndërveprueshmëria, shërbimet digjitale,
siguria kibernetike dhe inovacioni \cite{strategjia2023}. Në rrafshin
teknik, ekosistemi mbështetet në rrjetin e dedikuar qeveritar, në platformën
e ndërveprueshmërisë \textit{Government Gateway} (funksionale që nga viti
2017) dhe në portalin e-Kosova si dritare e vetme për shërbimet e të gjitha
niveleve të qeverisjes. Baza ligjore përfshin, ndër të tjera, Ligjin
Nr.~04/L-094 për shërbimet e shoqërisë së informacionit, Ligjin Nr.~08/L-022
për identifikimin elektronik dhe shërbimet e besuara, të përafruar me
eIDAS, dhe Ligjin Nr.~06/L-082 për mbrojtjen e të dhënave personale, të
përafruar me GDPR \cite{strategjia2023}. Procesi i integrimit evropian, i
formalizuar me Marrëveshjen e Stabilizim-Asociimit (2016) dhe me aplikimin
për anëtarësim në BE (2022) \cite{ec2024}, e shndërron përafrimin me
Kornizën Evropiane të Ndërveprueshmërisë (EIF) \cite{eif2017} nga një
zgjedhje politike në një detyrim strategjik.

Megjithatë, të dhënat e vetë dokumenteve strategjike zbulojnë një paradoks
ndërmjet lidhshmërisë dhe përdorimit (Tabela~\ref{tab:treguesit}). Ndonëse
pothuajse të gjitha ekonomitë familjare kanë qasje në internet, vetëm
5.3\% e përdoruesve të internetit kanë dorëzuar formularë online, ndërsa nga
mbi 700 shërbime që ofrojnë institucionet qendrore dhe lokale vetëm rreth
10\% ofrohen online \cite{strategjia2023,agjenda2030}. Sipas OECD-së, në vitin
2022 72\% e individëve nuk e kishin përdorur internetin për informata ose
shërbime të administratës publike \cite{oecd2024}. Në Indeksin e Maturitetit
GovTech (GTMI) të Bankës Botërore, Kosova ka kaluar nga grupi C në grupin B,
kryesisht falë ofrimit të shërbimeve, por mbetet më e dobët pikërisht te
mundësuesit institucionalë (GTEI) dhe angazhimi digjital i qytetarëve
\cite{strategjia2023}. Kufizimi kryesor, pra, nuk është infrastruktura, por
kapaciteti institucional për ta shndërruar atë në vlerë publike.

\begin{table}[!t]
\caption{Treguesit kryesorë të gjendjes dhe synimet e e-Qeverisjes në
Kosovë sipas dokumenteve strategjike}
\label{tab:treguesit}
\centering
\footnotesize
\renewcommand{\arraystretch}{1.2}
\begin{tabular}{>{\raggedright\arraybackslash}p{0.56\columnwidth}
                >{\raggedright\arraybackslash}p{0.34\columnwidth}}
\toprule
\textbf{Treguesi} & \textbf{Vlera / synimi} \\
\midrule
Ekonomi familjare me qasje në internet \cite{strategjia2023} &
99.8\% (2021--2022) \\
Përdorues të internetit që kanë dorëzuar formularë online
\cite{agjenda2030} & 5.3\% \\
Shërbime (nga mbi 700, qendrore e lokale) të ofruara online
\cite{strategjia2023} & $\approx$10\% \\
Shërbime qendrore e lokale të digjitalizuara -- synim
\cite{strategjia2023} & 40\% (2025); 65\% (2027) \\
Sisteme informative të lidhura me \textit{Government Gateway} -- synim
\cite{strategjia2023} & 30\% (2025); 70\% (2027) \\
Indeksi i mundësuesve GovTech (GTEI) \cite{strategjia2023} &
0.456, niveli C; synim $\geq$0.70 (2027) \\
Specialistë të TI-së në sektorin publik \cite{strategjia2023} &
$\approx$450, ose 0.5\% e punësimit publik (norma ndërkombëtare 2--3\%) \\
Buxheti indikativ i AD~2030 \cite{agjenda2030} &
152 mln euro, prej tyre 35\% nga donatorët \\
\bottomrule
\end{tabular}
\end{table}

Ky hendek shfaqet në tri dimensione të ndërlidhura. \textit{Horizontalisht},
shumë shërbime ofrohen përmes portaleve individuale të agjencive, metadatat
nuk janë të harmonizuara, nuk ekziston një pasqyrë e sistemeve të TI-së në
sektorin publik dhe arkitektura e ndërmarrjes së ``qeverisë në tërësi''
mungon ose është e vjetruar \cite{strategjia2023}; vlerësimi i Agjendës
2013--2020 tregoi gjithashtu se ministritë i zbatuan detyrat e tyre në
mënyra shumë të ndryshme, pa mekanizma korrigjues \cite{agjenda2030}.
\textit{Vertikalisht}, 38 komunat kanë kompetenca vetanake në ofrimin e
shërbimeve sipas Ligjit Nr.~03/L-040 për vetëqeverisjen lokale
\cite{ligji03L040}, por kapacitete digjitale shumë të pabarabarta, ndërkohë
që në përshkrimin e strukturave të zbatimit të Strategjisë 2023--2027 niveli
lokal nuk ka përfaqësim të dedikuar \cite{strategjia2023}.
\textit{Institucionalisht}, arkitektura e koordinimit
mbështetet në Komisionin për Transformim Digjital të
kryesuar nga Kryeministri, Njësinë e Transformimit
Digjital, komitetin teknik dhe Këshillin e RAP-it
\cite{strategjia2023}. Ndërkohë, kalimi i Agjencisë
së Shoqërisë së Informacionit (ASHI) në kuadër të
Ministrisë së Digjitalizimit dhe Administratës Publike
përfaqëson një hap të ri drejt konsolidimit institucional
dhe forcimit të koordinimit ndërinstitucional në
procesin e transformimit digjital; AD~2030 parashikon, nga ana e saj, raportim vjetor
nga Zyra e Kryeministrit dhe Ministria e Ekonomisë \cite{agjenda2030},
ndërsa në fushën e sigurisë kibernetike veprojnë akterë të tjerë, si
KOS-CERT, i cili sipas gjendjes së vitit 2022 funksiononte në kuadër të
ARKEP-it \cite{kipred2022}. Vetë regjistri i rreziqeve
të Strategjisë pranon se rolet e palëve janë të paqarta dhe se ka fusha të
mbivendosura, ndërsa menaxhimi i mundësuesve kryesorë teknikë (identiteti
digjital, \textit{Government Gateway}, e-Kosova dhe qendra e të dhënave)
varet nga partnerët e jashtëm \cite{strategjia2023}.

Situata rëndohet nga një deficit i matjes. Meqë Kosova nuk është shtet
anëtar i OKB-së, ajo nuk përfshihet në Indeksin e Zhvillimit të
e-Qeverisjes (EGDI) \cite{undesa2024} dhe ende nuk është pjesë e indeksit
DESI të BE-së \cite{agjenda2030}; njëkohësisht, shumë tregues të Strategjisë
2023--2027 kanë vlerë bazë ``N/A'' ose ``për t'u matur'' dhe trajtohen si
projekt-sugjerime deri në hartimin e Pasaportës së Treguesve
\cite{strategjia2023}. Studimet akademike për Kosovën janë kryesisht
juridike ose të përqendruara te pranimi i shërbimeve
\cite{dermaku2024,krasniqi2026} dhe rrallë e lidhin maturinë e shërbimeve me
kornizën ligjore e institucionale dhe me dinamikat qendër--komunë. Duke u
nisur nga premisa se sfida kryesore e e-Qeverisjes në Kosovë nuk është më
mungesa e teknologjisë, por mungesa e koordinimit institucional që e
shndërron teknologjinë në vlerë publike, ky punim propozon një vlerësim të
integruar të modelit aktual, të kornizës ligjore e institucionale, të
faktorëve SWOT në nivel qendror e lokal dhe të mekanizmave të koordinimit
ndërinstitucional.
% ---------------------------------------------------------------------
% 4. QËLLIMI I PUNIMIT
% ---------------------------------------------------------------------
\section{Qëllimi i punimit}
\label{sec:qellimi}
Qëllimi i këtij punimi është të vlerësojë në mënyrë kritike modelin aktual
të e-Qeverisjes në Kosovë dhe të identifikojë kushtet ligjore,
institucionale dhe organizative që e mundësojnë ose e pengojnë kalimin nga
digjitalizimi i shërbimeve të veçanta drejt një sistemi të integruar, të
sigurt dhe të orientuar te përdoruesi, në nivel qendror dhe lokal.
E-Qeverisja konceptualizohet si sistem socio-teknik, rezultatet e të cilit
varen nga ndërveprimi i politikave, menaxhimit dhe teknologjisë
\cite{pardo2012}. Kuadri analitik ndërthur modelin e fazave të
Layne dhe Lee \cite{layne2001} dhe modelin evolutiv për vlerësimin e maturisë, me katër shtresat e
ndërveprueshmërisë të EIF-it (ligjore, organizative, semantike dhe teknike)
për hartëzimin e kornizës normative dhe institucionale \cite{eif2017}.
Metodologjikisht, punimi mbështetet në analizë cilësore të dokumenteve
(AD~2030, Strategjia 2023--2027, legjislacioni përkatës, raportet
ndërkombëtare të monitorimit dhe studimet e politikave), të krahasuar me
treguesit sasiorë të disponueshëm; gjetjet sintetizohen përmes analizës SWOT
të diferencuar për nivelin qendror e lokal dhe zgjerohen në matricën TOWS,
në mënyrë që diagnoza të shndërrohet në opsione strategjike. Pyetja kryesore
kërkimore është: \textit{cilët faktorë ligjorë, institucionalë dhe
teknologjikë e shpjegojnë hendekun ndërmjet kapaciteteve digjitale të
ndërtuara dhe përdorimit real të shërbimeve elektronike në Kosovë, dhe cili
model i koordinimit ndërinstitucional mund ta ngushtojë këtë hendek ndërmjet
nivelit qendror dhe atij lokal?} Synimi përfundimtar është t'u ofrohet
politikëbërësve një bazë analitike dhe një udhërrëfyes i zbatueshëm, në
përputhje me afatet e Strategjisë 2023--2027 dhe të AD~2030.
% ---------------------------------------------------------------------
% 5. OBJEKTIVAT E PUNIMIT
% ---------------------------------------------------------------------
\section{Objektivat e punimit}
\label{sec:objektivat}
\begin{enumerate}
\renewcommand{\labelenumi}{O\arabic{enumi}:}
\item të vlerësohet niveli i maturisë së modelit aktual të e-Qeverisjes në
dimensionin horizontal (ndërsektorial) dhe vertikal (qendër--komunë) dhe të
formulohen rekomandime;
\item të hartëzohen boshllëqet, mbivendosjet e mandateve dhe mospërputhjet e
kornizës ligjore e institucionale sipas shtresave të EIF-it dhe
\textit{acquis}-së së BE-së;
\item të kryhet një analizë krahasuese SWOT për nivelin qendror dhe atë
lokal, e zgjeruar në matricë TOWS;
\item të propozohet një model qeverisjeje për koordinimin ndërinstitucional,
me role, instrumente dhe tregues të matshëm të performancës.
\end{enumerate}
% ---------------------------------------------------------------------
% 6. GJENDJA AKTUALE (STATE OF THE ART)
% ---------------------------------------------------------------------
\section{Gjendja aktuale e punëve të ngjashme (State of the Art)}
\label{sec:sota}
\subsection{Maturia dhe faktorët e zbatimit}
Literatura e hershme e konceptualizon e-Qeverisjen si proces me faza. Layne
dhe Lee \cite{layne2001} dallojnë katalogimin, transaksionin, integrimin
vertikal dhe integrimin horizontal, duke argumentuar se çdo fazë rrit
kompleksitetin teknologjik e organizativ, por edhe vlerën për qytetarin.
Lee e zhvendos fokusin nga teknologjia te
institucionet përmes një modeli evolutiv katërfazor (digjitalizim,
transformim, angazhim, kontekstualizim), ku faza më e avancuar kërkon
përshtatjen e politikave me nevojat specifike të vendit dhe të komunitetit.
Kalimi ndërmjet fazave nuk është linear: Glyptis et al. gjen se në
komuna e-Qeverisja mbetet kryesisht në fazat informative e transaksionale
dhe se madhësia e kapacitetet administrative ndikojnë ndjeshëm. Për vendet
e vogla, Glyptis et al. \cite{glyptis2020} e lidhin suksesin e projekteve me
gatishmërinë elektronike, kornizën politike e ligjore dhe faktorët
organizativë e socio-kulturorë.

\subsection{Ndërveprueshmëria dhe koordinimi}
Linja e dytë e trajton integrimin si problem qeverisjeje. Pardo, Nam dhe
Burke \cite{pardo2012} tregojnë empirikisht se ndërveprueshmëria është
rezultat i ndërveprimit të politikave, menaxhimit dhe teknologjisë, ku
qeverisja e përbashkët dhe besimi ndërinstitucional janë po aq të
rëndësishëm sa standardet teknike. Rasti i Estonisë, model referimi për
Strategjinë kosovare, tregon se suksesi lidhet më shumë me vendimet e hershme
arkitekturore, kulturën e inovacionit dhe bashkëpunimin publik--privat sesa
me vëllimin e investimeve \cite{kalvet2012}.

\subsection{Studimet për Kosovën dhe boshllëku kërkimor}
Dërmaku dhe Emini \cite{dermaku2024} analizojnë bazën juridike të
digjitalizimit të administratës dhe tregojnë se transformimi kërkon
përputhje ndërmjet teknologjisë dhe normave të procedurës administrative,
regjistrave publikë dhe mbrojtjes së të dhënave. Krasniqi, Elezi dhe Vuniqi
\cite{krasniqi2026} identifikojnë besimin institucional, aftësitë digjitale,
interoperabilitetin dhe sigurinë kibernetike si sfidat kryesore të
platformës e-Kosova. Në dimensionin e sigurisë, Peci dhe Ukshini
\cite{kipred2022} konstatojnë se korniza ligjore është ndërtuar gradualisht
që nga viti 2010, por se bashkëpunimi me akademinë është i herëpashershëm
dhe ai me qytetarët pothuajse mungon. Në sintezë, literatura ndërkombëtare ofron instrumente
të fuqishme për maturinë, faktorët e suksesit dhe ndërveprueshmërinë, por
këto rrallë integrohen në një kuadër të vetëm. Studimet për Kosovën janë
kryesisht sektoriale (juridike, të pranimit ose të sigurisë) dhe nuk e lidhin
vlerësimin e modelit me analizën SWOT qendër--komunë dhe me dizajnin e
mekanizmave të koordinimit. Ky punim synon ta adresojë pikërisht këtë
boshllëk.
% ---------------------------------------------------------------------
% 7. REFERENCAT (renditur sipas citimit të parë, stili IEEE)
% Punimet shkencore me DOI janë verifikuar përmes Crossref.
% ---------------------------------------------------------------------
\begin{thebibliography}{00}
\bibitem{strategjia2023} Qeveria e Republikës së Kosovës, \emph{Strategjia e Kosovës për Qeverisje Elektronike 2023--2027}, versioni final, Prishtinë, Kosovë, 2023.
\bibitem{agjenda2030} Ministria e Ekonomisë, \emph{Agjenda Digjitale e Kosovës 2030}, Prishtinë, Kosovë, 2023.
\bibitem{ec2024} European Commission, ``Kosovo 2024 Report,'' Commission Staff Working Document, Brussels, Belgium, 2024.
\bibitem{eif2017} European Commission, ``New European Interoperability Framework: Promoting seamless services and data flows for European public administrations,'' COM(2017) 134 final, Luxembourg: Publications Office of the European Union, 2017.
\bibitem{oecd2024} OECD, \emph{Western Balkans Competitiveness Outlook 2024: Kosovo}. Paris, France: OECD Publishing, 2024.
\bibitem{ligji03L040} Kuvendi i Republikës së Kosovës, ``Ligji Nr. 03/L-040 për Vetëqeverisjen Lokale,'' \emph{Gazeta Zyrtare e Republikës së Kosovës}, 2008.
\bibitem{kipred2022} L. Peci and V. Ukshini, \emph{Qeverisja e mirë në sigurinë kibernetike në Kosovë: Forcimi i themeleve dhe i institucioneve të reja}. Prishtinë, Kosovë: KIPRED/DCAF, Sep. 2022.
\bibitem{undesa2024} United Nations Department of Economic and Social Affairs, \emph{United Nations E-Government Survey 2024: Accelerating Digital Transformation for Sustainable Development}. New York, NY, USA: United Nations, 2024.
\bibitem{dermaku2024} K. Dërmaku and A. Emini, ``Digitisation of administration and legal basis in Kosovo,'' \emph{Access to Justice in Eastern Europe}, vol. 7, no. 1, pp. 245--263, 2024, doi: 10.33327/AJEE-18-7.1-a000107.
\bibitem{krasniqi2026} M. Krasniqi, Z. Elezi, and D. Vuniqi, ``The digitalization of public administration in Kosovo: Challenges, opportunities and the impact on the quality of public services,'' \emph{Prizren Social Science Journal}, vol. 10, no. 2, pp. 165--177, 2026, doi: 10.32936/pssj.v10i2.880.
\bibitem{pardo2012} T. A. Pardo, T. Nam, and G. B. Burke, ``E-government interoperability: Interaction of policy, management, and technology dimensions,'' \emph{Social Science Computer Review}, vol. 30, no. 1, pp. 7--23, 2012, doi: 10.1177/0894439310392184.
\bibitem{layne2001} K. Layne and J. Lee, ``Developing fully functional E-government: A four stage model,'' \emph{Government Information Quarterly}, vol. 18, no. 2, pp. 122--136, 2001, doi: 10.1016/S0740-624X(01)00066-1.
\bibitem{glyptis2020} L. Glyptis, M. Christofi, D. Vrontis, M. Del Giudice, S. Dimitriou, and P. Michael, ``E-Government implementation challenges in small countries: The project manager's perspective,'' \emph{Technological Forecasting and Social Change}, vol. 152, 119880, 2020, doi: 10.1016/j.techfore.2019.119880.
\bibitem{kalvet2012} T. Kalvet, ``Innovation: A factor explaining e-government success in Estonia,'' \emph{Electronic Government, an International Journal}, vol. 9, no. 2, pp. 142--157, 2012, doi: 10.1504/EG.2012.046266.
\end{thebibliography}
\end{document}
